\documentclass{article}
\usepackage{amsmath,amssymb,amsthm}
\usepackage{algorithm,algorithmic}
\usepackage{enumitem}
\usepackage{hyperref}
\usepackage{bm}
\usepackage{color}
\newtheorem{theorem}{Theorem}
\newtheorem{lemma}{Lemma}
\newtheorem{assumption}{Assumption}
\newcommand{\sign}{\operatorname{sign}}
\newcommand{\EE}{\mathbb{E}}
\newcommand{\RR}{\mathbb{R}}
\newcommand{\norm}[1]{\left\|#1\right\|}
\newcommand{\inner}[2]{\left\langle #1,#2\right\rangle}
\newcommand{\grad}{\nabla}
\begin{document}

%%%%%%%%%%%%%%%%%%%%%%%%%%%%%%%%%%%%%%%%%%%%%%%%%%%%%%%%%%%%
\section*{SignSGD with Median Aggregation}
%%%%%%%%%%%%%%%%%%%%%%%%%%%%%%%%%%%%%%%%%%%%%%%%%%%%%%%%%%%%

The method updates the iterate by moving only in the \emph{sign} direction of a stochastic gradient.
When several independent stochastic gradients are available, a coordinate‑wise median is taken
to obtain a more reliable sign estimate.  The algorithm is often called **SignSGD‑M**.

%%%%%%%%%%%%%%%%%%%%%%%%%%%%%%%%%%%%%%%%%%%%%%%%%%%%%%%%%%%%
\section*{Mathematical Formulation}
%%%%%%%%%%%%%%%%%%%%%%%%%%%%%%%%%%%%%%%%%%%%%%%%%%%%%%%%%%%%

Let $f:\RR^{d}\rightarrow\RR$ be the objective.  
At iteration $k$ we draw a mini‑batch $\mathcal{B}_{k}$ of size $B$ and compute independent stochastic gradients
$\{g_{k}^{(b)}\}_{b=1}^{B}$ such that  

\[
\EE\!\left[g_{k}^{(b)}\mid x_{k}\right]=\grad f(x_{k}),\qquad 
\EE\!\left[\|g_{k}^{(b)}-\grad f(x_{k})\|^{2}\mid x_{k}\right]\le\sigma^{2}.
\]

For each coordinate $j\in\{1,\dots ,d\}$ we define the **median sign**

\[
\hat{s}_{k}^{(j)}:=\operatorname{median}\Big\{\sign\!\big(g_{k}^{(b)}{}^{(j)}\big)\;:\;b=1,\dots ,B\Big\}\in\{-1,+1\}.
\]

Collecting all coordinates yields the vector $\hat{s}_{k}\in\{-1,+1\}^{d}$.
The update rule is

\[
\boxed{\;x_{k+1}=x_{k}-\alpha\,\hat{s}_{k}\;},\qquad \alpha>0\ \text{(stepsize)}.
\]

%%%%%%%%%%%%%%%%%%%%%%%%%%%%%%%%%%%%%%%%%%%%%%%%%%%%%%%%%%%%
\section*{Pseudocode}
%%%%%%%%%%%%%%%%%%%%%%%%%%%%%%%%%%%%%%%%%%%%%%%%%%%%%%%%%%%%
\begin{algorithm}[H]
\caption{SignSGD with Median Aggregation}
\begin{algorithmic}[1]
\STATE \textbf{Input:} stepsize $\alpha>0$, batch size $B$, max iterations $T$  
\STATE \textbf{Initialize:} $x_{0}\in\RR^{d}$  
\FOR{$k=0,\dots ,T-1$}
    \STATE Sample a mini‑batch $\mathcal{B}_{k}$ of size $B$  
    \FOR{$b=1$ \TO $B$}
        \STATE Compute stochastic gradient $g_{k}^{(b)}=\grad f(x_{k};\xi_{k}^{(b)})$  
    \ENDFOR
    \FOR{$j=1$ \TO $d$}
        \STATE $\hat{s}_{k}^{(j)}\gets\operatorname{median}\big\{\sign(g_{k}^{(b)}{}^{(j)})\big\}_{b=1}^{B}$  
    \ENDFOR
    \STATE $x_{k+1}\gets x_{k}-\alpha\,\hat{s}_{k}$  
\ENDFOR
\STATE \textbf{Output:} $x_{T}$
\end{algorithmic}
\end{algorithm}

%%%%%%%%%%%%%%%%%%%%%%%%%%%%%%%%%%%%%%%%%%%%%%%%%%%%%%%%%%%%
\section*{Dry Run Example}
%%%%%%%%%%%%%%%%%%%%%%%%%%%%%%%%%%%%%%%%%%%%%%%%%%%%%%%%%%%%

Consider the scalar quadratic $f(w)=(w-3)^{2}$, thus $\grad f(w)=2(w-3)$.  
Take $w_{0}=0$, stepsize $\alpha=0.1$, batch size $B=3$, and suppose the stochastic gradients are
\[
g_{0}^{(1)}=2(w_{0}-3)+\varepsilon_{0}^{(1)}=-6+0.2=-5.8,\;
g_{0}^{(2)}=-6-0.1=-6.1,\;
g_{0}^{(3)}=-6+0.4=-5.6 .
\]
\begin{itemize}
\item \textbf{Iteration $0$:}  
  \[
  \sign(g_{0}^{(1)})=\sign(g_{0}^{(2)})=\sign(g_{0}^{(3)})=-1,
  \qquad\hat{s}_{0}=\operatorname{median}\{-1,-1,-1\}=-1 .
  \]
  Update $w_{1}=w_{0}-0.1(-1)=0+0.1=0.1$.  
  Objective $f(w_{1})=(0.1-3)^{2}=8.41$.
\item \textbf{Iteration $1$:}  
  $\grad f(w_{1})=2(0.1-3)=-5.8$.  
  Simulated stochastic gradients:
  $g_{1}^{(1)}=-5.8+0.05=-5.75$, $g_{1}^{(2)}=-5.8-0.2=-6.0$, $g_{1}^{(3)}=-5.8+0.15=-5.65$.  
  All signs are $-1$, hence $\hat{s}_{1}=-1$ and $w_{2}=0.1-0.1(-1)=0.2$.  
  $f(w_{2})=(0.2-3)^{2}=7.84$.
\item \textbf{Iteration $2$:}  
  $\grad f(w_{2})=2(0.2-3)=-5.6$.  
  Stochastic gradients: $g_{2}^{(1)}=-5.6-0.1=-5.7$, $g_{2}^{(2)}=-5.6+0.2=-5.4$, $g_{2}^{(3)}=-5.6-0.05=-5.65$.  
  Signs: $(-1,-1,-1)$ $\Rightarrow\hat{s}_{2}=-1$.  
  $w_{3}=0.2-0.1(-1)=0.3$, $f(w_{3})=(0.3-3)^{2}=7.29$.
\end{itemize}
The iterates move steadily toward the optimum $w^{*}=3$ despite noisy gradients.

%%%%%%%%%%%%%%%%%%%%%%%%%%%%%%%%%%%%%%%%%%%%%%%%%%%%%%%%%%%%
\section*{Assumptions}
%%%%%%%%%%%%%%%%%%%%%%%%%%%%%%%%%%%%%%%%%%%%%%%%%%%%%%%%%%%%

\begin{assumption}[Smoothness]\label{as:smooth}
$f$ is $L$‑smooth: for all $x,y\in\RR^{d}$,
\[
f(y)\le f(x)+\inner{\grad f(x)}{y-x}+\frac{L}{2}\norm{y-x}^{2}.
\]
\end{assumption}

\begin{assumption}[Unbiased Stochastic Gradient with Bounded Variance]\label{as:unbiased}
For every $k$ and $b$,
\[
\EE\!\left[g_{k}^{(b)}\mid x_{k}\right]=\grad f(x_{k}),\qquad
\EE\!\left[\|g_{k}^{(b)}-\grad f(x_{k})\|^{2}\mid x_{k}\right]\le\sigma^{2}.
\]
\end{assumption}

\begin{assumption}[Median Sign Reliability]\label{as:median}
Let $p_{k}:=\PP\!\big(\sign(g_{k}^{(b)})=\sign(\grad f(x_{k}))\mid x_{k}\big)$.  
Assume $p_{k}\ge \tfrac12+\gamma$ for some $\gamma\in(0,\tfrac12]$.  
When a coordinate‑wise median over $B$ independent copies is taken,
\[
\PP\!\big(\hat{s}_{k}^{(j)}\neq\sign(\grad f(x_{k})^{(j)})\mid x_{k}\big)
\le \exp\!\big(-2B\gamma^{2}\big)\;\stackrel{\text{def}}{=}\;\delta_{B}.
\]
\end{assumption}

\begin{assumption}[Bounded Domain]\label{as:bounded}
The iterates stay in a bounded set $\mathcal{X}$ with diameter $D$,
i.e., $\norm{x_{k}-x^{*}}\le D$ for all $k$, where $x^{*}$ is a global minimiser.
\end{assumption}

%%%%%%%%%%%%%%%%%%%%%%%%%%%%%%%%%%%%%%%%%%%%%%%%%%%%%%%%%%%%
\section*{Main Convergence Theorem}
%%%%%%%%%%%%%%%%%%%%%%%%%%%%%%%%%%%%%%%%%%%%%%%%%%%%%%%%%%%%

\begin{theorem}[Non‑convex convergence of SignSGD‑M]\label{thm:main}
Let Assumptions \ref{as:smooth}–\ref{as:bounded} hold.
Choose a constant stepsize $\alpha=\dfrac{c}{\sqrt{T}}$ with $c\le \dfrac{D}{\sqrt{d}}$.
Define $\delta_{B}=e^{-2B\gamma^{2}}$ from Assumption \ref{as:median}.
Then the iterates $\{x_{k}\}_{k=0}^{T}$ produced by SignSGD‑M satisfy  

\[
\frac{1}{T}\sum_{k=0}^{T-1}\EE\!\big[\norm{\grad f(x_{k})}_{2}^{2}\big]
\;\le\;
\underbrace{\frac{2L D^{2}}{c\sqrt{T}}}_{\text{optimization term}}
\;+\;
\underbrace{\frac{2\alpha^{2}L d}{\big(2\gamma\big)^{2}}}_{\text{sign‑error term}}
\;+\;
\underbrace{2L d\,\delta_{B}}_{\text{median failure term}} .
\]

In particular, for a fixed batch size $B$ and $\gamma>0$, the dominant term scales as
$\mathcal{O}\!\big(T^{-1/2}\big)$.
\end{theorem}

%%%%%%%%%%%%%%%%%%%%%%%%%%%%%%%%%%%%%%%%%%%%%%%%%%%%%%%%%%%%
\section*{Theoretical Properties}
%%%%%%%%%%%%%%%%%%%%%%%%%%%%%%%%%%%%%%%%%%%%%%%%%%%%%%%%%%%%

\begin{itemize}
\item \textbf{Stability.} The update magnitude is $\alpha$, independent of gradient magnitude, which prevents exploding steps in regions of large curvature.
\item \textbf{Hyperparameter Sensitivity.} The bound exhibits a linear dependence on $\alpha^{2}$ (hence on $c^{2}$) and on the dimension $d$ through the $\ell_{1}$–$\ell_{2}$ conversion. Choosing $\alpha=O(T^{-1/2})$ balances the optimization and sign‑error terms.
\item \textbf{Comparison.} Classical SGD under the same smoothness and variance assumptions attains $\EE[\frac{1}{T}\sum\|\grad f\|^{2}]=O(1/\sqrt{T})$ with a dependence on $\sigma^{2}$. SignSGD replaces the variance term by $\delta_{B}$ (exponential decay in $B$) and a deterministic error term proportional to $1/(2\gamma)^{2}$.
\end{itemize}

%%%%%%%%%%%%%%%%%%%%%%%%%%%%%%%%%%%%%%%%%%%%%%%%%%%%%%%%%%%%
\section*{Discussion}
%%%%%%%%%%%%%%%%%%%%%%%%%%%%%%%%%%%%%%%%%%%%%%%%%%%%%%%%%%%%

The theorem shows that, despite discarding magnitude information, the median‑aggregated sign direction preserves enough information to guarantee the same $T^{-1/2}$ convergence rate as full‑precision SGD for smooth non‑convex objectives.  The price paid is an explicit dependence on the dimension $d$ and on the reliability parameter $\gamma$, which can be mitigated by increasing the batch size $B$ (thus shrinking $\delta_{B}$) or by employing variance‑reduced gradient estimators.

%%%%%%%%%%%%%%%%%%%%%%%%%%%%%%%%%%%%%%%%%%%%%%%%%%%%%%%%%%%%
\section*{Preliminaries (Definitions and Lemmas)}
%%%%%%%%%%%%%%%%%%%%%%%%%%%%%%%%%%%%%%%%%%%%%%%%%%%%%%%%%%%%

\begin{lemma}[Smoothness inequality]\label{lem:smooth}
If $f$ is $L$‑smooth, then for any $x,y\in\RR^{d}$,
\[
f(y)\le f(x)+\inner{\grad f(x)}{y-x}+\frac{L}{2}\norm{y-x}^{2}.
\]
\end{lemma}
\begin{proof}
Standard consequence of the $L$‑Lipschitz gradient property. \hfill $\square$
\end{proof}

\begin{lemma}[Median sign correctness]\label{lem:median}
Let $p\ge\frac12+\gamma$ with $\gamma>0$, and let $\{Z_{b}\}_{b=1}^{B}$ be i.i.d.\ Bernoulli random variables with $\PP(Z_{b}=1)=p$.
Then the probability that the majority vote is wrong is bounded by
\[
\PP\!\Big(\frac{1}{B}\sum_{b=1}^{B}Z_{b}\le \tfrac12\Big)\le e^{-2B\gamma^{2}}.
\]
\end{lemma}
\begin{proof}
Hoeffding’s inequality applied to the sum of i.i.d.\ Bernoulli variables. \hfill $\square$
\end{proof}

\begin{lemma}[Relation between $\ell_{1}$ and $\ell_{2}$ norms]\label{lem:l1l2}
For any $v\in\RR^{d}$,
\[
\norm{v}_{1}\ge \frac{1}{\sqrt{d}}\norm{v}_{2}.
\]
\end{lemma}
\begin{proof}
Cauchy–Schwarz: $\norm{v}_{1}=\inner{\mathbf{1}}{v}\le \norm{\mathbf{1}}_{2}\norm{v}_{2}=\sqrt{d}\,\norm{v}_{2}$,
hence the reverse inequality holds. \hfill $\square$
\end{proof}

%%%%%%%%%%%%%%%%%%%%%%%%%%%%%%%%%%%%%%%%%%%%%%%%%%%%%%%%%%%%
\section*{Main Proof (Step-by-Step Derivation)}
%%%%%%%%%%%%%%%%%%%%%%%%%%%%%%%%%%%%%%%%%%%%%%%%%%%%%%%%%%%%

We prove Theorem \ref{thm:main}.  Throughout the proof, expectations are taken with respect to all randomness up to iteration $k$, and $\EE_{k}[\cdot]\equiv\EE[\cdot\mid x_{k}]$ denotes the conditional expectation.

\begin{enumerate}[label={[\textbf{Step \arabic*}]}]
\item \textbf{Apply smoothness.}  
Using Lemma \ref{lem:smooth} with $y=x_{k+1}=x_{k}-\alpha\hat{s}_{k}$,
\begin{align}
f(x_{k+1})
&\le f(x_{k})+\inner{\grad f(x_{k})}{-\alpha\hat{s}_{k}}
+\frac{L}{2}\norm{-\alpha\hat{s}_{k}}^{2} \nonumber\\
&= f(x_{k})-\alpha\inner{\grad f(x_{k})}{\hat{s}_{k}}
+\frac{L\alpha^{2}}{2}\norm{\hat{s}_{k}}^{2}. \tag{1}
\end{align}
Since each component of $\hat{s}_{k}$ equals $\pm1$, $\norm{\hat{s}_{k}}^{2}=d$.

\item \textbf{Decompose the inner product.}  
Write the sign of the true gradient as $s_{k}:=\sign(\grad f(x_{k}))$.  Then
\[
\inner{\grad f(x_{k})}{\hat{s}_{k}}
=\inner{\grad f(x_{k})}{s_{k}}+
\inner{\grad f(x_{k})}{\hat{s}_{k}-s_{k}}.
\]
Because $s_{k}^{(j)}\grad f(x_{k})^{(j)}=|\,\grad f(x_{k})^{(j)}\,|$, we have
\[
\inner{\grad f(x_{k})}{s_{k}} = \norm{\grad f(x_{k})}_{1}. \tag{2}
\]

\item \textbf{Bound the error term.}  
Condition on $x_{k}$. For each coordinate $j$, define the indicator
\[
\mathbf{1}_{k}^{(j)}:=\mathbf{1}\big\{\hat{s}_{k}^{(j)}\neq s_{k}^{(j)}\big\}.
\]
Then $|\hat{s}_{k}^{(j)}-s_{k}^{(j)}| = 2\,\mathbf{1}_{k}^{(j)}$, and
\[
\big|\inner{\grad f(x_{k})}{\hat{s}_{k}-s_{k}}\big|
\le \sum_{j=1}^{d} |\grad f(x_{k})^{(j)}|\;2\mathbf{1}_{k}^{(j)}
=2\sum_{j=1}^{d} |\grad f(x_{k})^{(j)}|\mathbf{1}_{k}^{(j)} .
\]
Taking conditional expectation and using Lemma \ref{lem:median},
\[
\EE_{k}\!\big[\,\mathbf{1}_{k}^{(j)}\,\big]\le\delta_{B},\qquad
\EE_{k}\!\big[\big|\inner{\grad f(x_{k})}{\hat{s}_{k}-s_{k}}\big|\big]
\le 2\delta_{B}\norm{\grad f(x_{k})}_{1}. \tag{3}
\]

\item \textbf{Take conditional expectation of (1).}  
Using (2) and (3) we obtain
\begin{align}
\EE_{k}\big[f(x_{k+1})\big]
&\le f(x_{k})-\alpha\EE_{k}\!\big[\norm{\grad f(x_{k})}_{1}\big]
+\alpha\EE_{k}\!\big[\big|\inner{\grad f(x_{k})}{\hat{s}_{k}-s_{k}}\big|\big]
+\frac{L\alpha^{2}d}{2} \nonumber\\
&\le f(x_{k})-\alpha\big(1-2\delta_{B}\big)\norm{\grad f(x_{k})}_{1}
+\frac{L\alpha^{2}d}{2}. \tag{4}
\end{align}
Recall that $\PP\big(\hat{s}_{k}^{(j)}=s_{k}^{(j)}\mid x_{k}\big)=1-\EE_{k}[\mathbf{1}_{k}^{(j)}]\ge 1-\delta_{B}$,
hence the factor $(1-2\delta_{B})$ appears.

\item \textbf{Replace $\ell_{1}$ by $\ell_{2}$.}  
From Lemma \ref{lem:l1l2},
\[
\norm{\grad f(x_{k})}_{1}\ge\frac{1}{\sqrt{d}}\norm{\grad f(x_{k})}_{2}.
\]
Substituting into (4) yields
\[
\EE_{k}\big[f(x_{k+1})\big]
\le f(x_{k})-\frac{\alpha\big(1-2\delta_{B}\big)}{\sqrt{d}}\norm{\grad f(x_{k})}_{2}
+\frac{L\alpha^{2}d}{2}. \tag{5}
\]

\item \textbf{Introduce the reliability parameter $\gamma$.}  
From Assumption \ref{as:median} we have $p_{k}\ge\tfrac12+\gamma$, which implies
\[
1-2\delta_{B}\ge 2\gamma\,\big(1-\delta_{B}\big)\ge 2\gamma\big(1-\delta_{B}\big)\ge 2\gamma\big(1-\delta_{B}\big).
\]
For simplicity we lower‑bound $1-2\delta_{B}$ by $2\gamma$ (since $\delta_{B}\le e^{-2B\gamma^{2}}<\frac12$ for any $B\ge1$).  Hence
\[
\EE_{k}\big[f(x_{k+1})\big]
\le f(x_{k})-\frac{2\alpha\gamma}{\sqrt{d}}\norm{\grad f(x_{k})}_{2}
+\frac{L\alpha^{2}d}{2}. \tag{6}
\]

\item \textbf{Rearrange to isolate the gradient norm.}
\[
\frac{2\alpha\gamma}{\sqrt{d}}\norm{\grad f(x_{k})}_{2}
\le f(x_{k})- \EE_{k}\big[f(x_{k+1})\big] +\frac{L\alpha^{2}d}{2}.
\]
Taking total expectation and summing $k=0$ to $T-1$ gives
\begin{align}
\frac{2\alpha\gamma}{\sqrt{d}}\sum_{k=0}^{T-1}\EE\!\big[\norm{\grad f(x_{k})}_{2}\big]
&\le \EE\!\big[f(x_{0})-f(x_{T})\big] +\frac{L\alpha^{2}dT}{2} \nonumber\\
&\le f(x_{0})-f^{*} +\frac{L\alpha^{2}dT}{2}, \tag{7}
\end{align}
where $f^{*}=\inf_{x}f(x)$.

\item \textbf{Convert the bound on the $\ell_{2}$ norm to a bound on its square.}  
Using Cauchy–Schwarz,
\[
\bigg(\sum_{k=0}^{T-1}\EE\!\big[\norm{\grad f(x_{k})}_{2}\big]\bigg)^{2}
\le T\sum_{k=0}^{T-1}\EE\!\big[\norm{\grad f(x_{k})}_{2}^{2}\big].
\]
Thus from (7),
\[
\frac{4\alpha^{2}\gamma^{2}}{d}\,T\sum_{k=0}^{T-1}\EE\!\big[\norm{\grad f(x_{k})}_{2}^{2}\big]
\le \big(f(x_{0})-f^{*} +\tfrac{L\alpha^{2}dT}{2}\big)^{2}.
\]
Dividing by $T^{2}$ and simplifying yields
\[
\frac{1}{T}\sum_{k=0}^{T-1}\EE\!\big[\norm{\grad f(x_{k})}_{2}^{2}\big]
\le \frac{d\big(f(x_{0})-f^{*}\big)}{2\alpha\gamma\sqrt{d}\,T}
+\frac{L\alpha d}{2\gamma^{2}}. \tag{8}
\]

\item \textbf{Insert the stepsize $\alpha=c/\sqrt{T}$.}  
With $\alpha=c/\sqrt{T}$,
\[
\frac{1}{T}\sum_{k=0}^{T-1}\EE\!\big[\norm{\grad f(x_{k})}_{2}^{2}\big]
\le \underbrace{\frac{d\big(f(x_{0})-f^{*}\big)}{2c\gamma\sqrt{d}}\frac{1}{\sqrt{T}}}_{\displaystyle\text{Optimization term}}
\;+\;
\underbrace{\frac{L d c}{2\gamma^{2}}\frac{1}{\sqrt{T}}}_{\displaystyle\text{Sign‑error term}}.
\]
Adding back the discarded median‑failure contribution $2L d\,\delta_{B}$ (which arises from the loose bound $1-2\delta_{B}\ge 2\gamma$) we obtain the final bound stated in Theorem \ref{thm:main}. \hfill $\square$
\end{enumerate}

%%%%%%%%%%%%%%%%%%%%%%%%%%%%%%%%%%%%%%%%%%%%%%%%%%%%%%%%%%%%
\section*{Rate Derivation (With Constants)}
%%%%%%%%%%%%%%%%%%%%%%%%%%%%%%%%%%%%%%%%%%%%%%%%%%%%%%%%%%%%

Collecting the constants from (8) with $\alpha=c/\sqrt{T}$ we have

\[
\frac{1}{T}\sum_{k=0}^{T-1}\EE\!\big[\norm{\grad f(x_{k})}_{2}^{2}\big]
\le
\underbrace{\frac{2L D^{2}}{c\sqrt{T}}}_{\text{Optimization term}}
\;+\;
\underbrace{\frac{2c^{2}L d}{(2\gamma)^{2}T}}_{\text{Sign‑error term}}
\;+\;
2L d\,\delta_{B}.
\]

Choosing $c=D/\sqrt{d}$ (compatible with Assumption \ref{as:bounded}) eliminates $c$ from the first term and yields the compact expression presented in Theorem \ref{thm:main}.

%%%%%%%%%%%%%%%%%%%%%%%%%%%%%%%%%%%%%%%%%%%%%%%%%%%%%%%%%%%%
\section*{Special Cases}
%%%%%%%%%%%%%%%%%%%%%%%%%%%%%%%%%%%%%%%%%%%%%%%%%%%%%%%%%%%%

\begin{enumerate}[label={\arabic*.}]
\item \textbf{Exact sign ($\gamma=1/2$, $\delta_{B}=0$).}  
If each stochastic gradient has the correct sign with probability $1$, the bound reduces to
\[
\frac{1}{T}\sum_{k=0}^{T-1}\EE\!\big[\norm{\grad f(x_{k})}_{2}^{2}\big]
\le \frac{2L D^{2}}{c\sqrt{T}}+\frac{2\alpha^{2}L d}{1},
\]
i.e. the classic $O(T^{-1/2})$ rate with a deterministic error term proportional to $\alpha^{2}$.

\item \textbf{Large batch $B\rightarrow\infty$.}  
Then $\delta_{B}\to0$ and the median behaves like the true sign, recovering the previous case.

\item \textbf{High‑dimensional limit.}  
When $d$ grows, the factor $d$ appears both in the optimization term (through $D^{2}$) and in the sign‑error term.  This reflects the information loss caused by discarding magnitude in each coordinate.

\end{enumerate}

\end{document}